\chapter{Diskussion und Fazit}
\label{kap:diskussion}

\section{Was ist belegt, was ist nur begründet?}
\label{sec:tragend}

Kapitel~\ref{kap:pipeline} hat gezeigt, dass die Pipeline die Anforderungen A1
--A6 erfüllt. Das ist eine Aussage über den Entwurf, nicht über die Wirkung.
Für eine belastbare Bilanz ist zwischen Entscheidungen zu unterscheiden, deren
Nutzen \emph{gemessen} wurde, und solchen, die lediglich \emph{begründet} sind.
Von stark nach schwach geordnet ergeben sich vier Stufen.

\textbf{Abliert}, also in eigenen Trainingsläufen geprüft, die sich nur in
dieser einen Größe unterscheiden, ist allein das Bündel der
Einbettungsmaßnahmen gegen S1--S5 samt Schlagschatten
(Abschnitt~\ref{sec:ablationC}); welche Einzelmaßnahme darin wie viel trägt,
ist offen. \textbf{Ohne Neutraining gemessen} sind die Inferenzparameter
(Abschnitt~\ref{sec:imgsz}). Nur die Auflösung zeigt dabei einen Effekt über
dem Bootstrap-Intervall; NMS-Schwelle und Test-Time-Augmentation bleiben beim
Ende-zu-Ende-Kopf wirkungslos, und \code{max\_det} verändert ab 100 den
Betriebspunkt und ab 200 auch \mAPfifty{} nicht mehr.
\textbf{An der erzeugten Verteilung kalibriert} sind Objektgrößenbereich,
Ovalisierungsanteil, Gitterabstände, die Basisgröße je Cluster, das
Tiefenmodell 75/25 und die Randregel: Die erzeugte Statistik trifft die reale
(Tabelle~\ref{tab:kalibrierung}), ihr Nutzen für die Detektion ist damit aber
nicht gemessen. \textbf{Nur begründet} sind alle übrigen Entscheidungen,
darunter die Herkunft der Okkluder, die Matting-Schwelle, Gamma, HSV, Mulde,
globale Bildstufe, Negativbildanteil sowie Dellen- und Perspektivanteil.

\absatz{Bilanz der Begründungen.} Die Zielsetzung hatte drei zulässige Quellen
festgelegt und für alles andere die Kennzeichnung als Annahme verlangt.
Tabelle~\ref{tab:parameter} dokumentiert 37 Transformationen mit 44
Parameterwerten: 7 sind an den Realdaten gemessen, 8 literaturgestützt, 2
physikalisch begründet, 11 gemischt und \textbf{16 ausdrücklich als Annahme
gekennzeichnet}. Rund 64\,\% sind damit auf eine der drei zulässigen Quellen
zurückgeführt -- die elf gemischten allerdings nur zur Hälfte --, 36\,\% sind
Annahmen. Diese Bilanz ist nicht nebensächlich: Bei rein synthetischen Trainingsdaten
ist die Liste dieser 44 Werte die vollständige Beschreibung der
Trainingsverteilung.

Gemessen an diesen 44 Werten ist ein einziges abliertes Bündel wenig. Die
Begründungen tragen den Entwurf, experimentell belegt ist daraus ein
Punkt. Er ist allerdings nicht beliebig gewählt, denn er betrifft das
Leitprinzip der Pipeline.

\absatz{Eine Größe verbindet fünf Entscheidungen.} Rückblickend steht eine
einzige Kenngröße hinter dem größten Teil der Arbeit: die effektive Objektgröße
$d_{\text{eff}} = d \cdot M/N$ aus Gleichung~\eqref{eq:deff}, bezogen auf den
Stride $S$ der feinsten Detektionsebene. Aus diesem Verhältnis folgen
\emph{erstens} die Kalibrierung des Generators auf 12--44\px, weil die
Realobjekte 6--36\px{} groß sind (Abschnitt~\ref{sec:valtest}),
\emph{zweitens} die Trainingsauflösung 1024, weil sie der Leinwandgröße
entspricht und die kalibrierte Verteilung damit unskaliert beim Netz ankommt
(Abschnitt~\ref{sec:trainingsaufbau}), \emph{drittens} die Wahl des P2-Kopfes,
weil er $S$ von 8 auf 4 senkt und die kritische Größe $S\sqrt{2}$ damit unter
alle auftretenden Objektgrößen schiebt (Abschnitt~\ref{sec:trainingsaufbau}),
\emph{viertens} die Inferenzauflösung (Abschnitt~\ref{sec:imgsz}) und
\emph{fünftens} die Erklärung des größten gemessenen Einzeleffekts, des
Einbruchs auf \mAPfifty{} \zswValNiedrig{} bei \code{imgsz}~640, den
Abschnitt~\ref{subsec:kleine-objekte} vorhergesagt hatte, bevor er gemessen
wurde. Diese Kette ist das methodisch belastbarste Ergebnis der Arbeit: Aus
einer einzigen Größe folgen fünf Entscheidungen, und eine daraus abgeleitete
Vorhersage ist eingetroffen.

\section{Grenzen und nächste Schritte}
\label{sec:grenzen}

\absatz{Grenzen der Pipeline.} \emph{(1)} Nur sechs Quellobjekte aus einer
einzigen Aufnahme -- jedes erscheint rund 22\,700-mal als gelabeltes Objekt.
Die photometrische und geometrische Streuung mildert das, hebt es aber nicht
auf; Augmentierung erzeugt keine neue Information über die Objektklasse. Keines
der sechs zeigt eine mit Laub oder Rost gefüllte Bohrung, und genau solche
Enden übersieht das Modell auf \code{val} (Abschnitt~\ref{sec:bootstrap}). Auf
die geringe Formvielfalt stützt sich auch die Wahl der Modellgröße
\enquote{s} (Abschnitt~\ref{sec:trainingsaufbau}); dass mehr Kapazität hier vor
allem auswendig lernen würde, ist plausibel \parencite{zhang2017rethinking},
aber nicht gemessen. \emph{(2)} Nur diffuses Licht im Quellmaterial: Harte
Schlagschatten, Gegenlicht, Nässe und Schnee fehlen vollständig und werden nur
approximativ über die globale photometrische Stufe angenähert. Ob das die
Detektion beeinträchtigt, zeigt die Fehleranalyse nicht; ihr zweiter Fehlertyp
geht auf die Erscheinung der Rohrenden selbst zurück. \emph{(3)} Reflect-Padding erzeugt Spiegelsymmetrie, wenn ein
Hintergrund für den gewünschten Zoom zu klein ist; eine spiegelsymmetrische
Bildhälfte ist selbst ein erkennbares Merkmal, die Maßnahme gegen Signatur~S1
erzeugt damit eine schwächere neue Auffälligkeit. \emph{(4)} Die Schaftrichtung ist an der
Säulenrichtung des Clusters kalibriert, nicht an der Rohrrichtung des
Hintergrunds. In 19,4\,\% der Cluster auf Clean Plates kreuzt der Schaft die
Hintergrundrohre deshalb unter mehr als $60^\circ$, betroffen sind 22,0\,\% der
erzeugten Bilder (\code{src/schaftwinkel.py}), obwohl real alle Rohre eines
Stapels parallel liegen. Ein Angleich wäre
mit geringem Aufwand möglich, weil die dominante Hintergrundrichtung gemessen
vorliegt (\code{src/occluders.py}).

\absatz{Grenzen der Bewertung.} Die schwerwiegendste methodische Schwäche liegt
nicht in der Pipeline, sondern in der Bewertung: \emph{je ein Bild} pro Split.
Daraus folgen breite Konfidenzintervalle: $[\zbootValLo;\,\zbootValHi]$ auf
\code{val} und $[\zbootTestLo;\,\zbootTestHi]$ auf \code{test}, also drei bis
vier Prozentpunkte Halbbreite -- Unterschiede unterhalb dieser Größe sind nicht
entscheidbar; beide Bilder zeigen dieselbe Objektart bei ähnlichem
Wetter, die Ergebnisse sagen also nichts über andere Zaunhersteller,
Rohrdurchmesser oder Lichtsituationen. Schließlich wird \code{val} doppelt
verwendet -- für die Modellauswahl \emph{und} für die Wahl der
Inferenzauflösung. \code{val} ist damit kein unabhängiger Bewertungssplit
mehr.

\absatz{Quellmaterial und Bewertungsdaten.} Bildinhalt teilen Quellmaterial und
Bewertungsbilder nicht (Abschnitt~\ref{sec:valtest}), unabhängig sind sie aber
nicht. Alle sechs Objektausschnitte und zwei der vier Clean Plates stammen aus
derselben Aufnahmeserie wie \code{val}, aufgenommen innerhalb von 20 Sekunden
am selben Ort mit derselben Kamera. Die \code{val}-Zahlen messen deshalb auch,
wie gut das Modell auf eine Szene überträgt, aus der sein Quellmaterial stammt.
Außerdem ist die Labelstatistik beider Bilder in die Kalibrierung eingegangen,
am deutlichsten bei der unteren Größengrenze von 12\px, dem 5\,\%-Quantil von
\code{test} (Befund~5). Streng getrennt wäre eine Kalibrierung allein an
\code{val}; die untere Grenze läge dann bei 19\px, und die kleinen Objekte in
\code{test} deckte nur noch die Trainingsskalierung ab. \code{test} ist damit
der einzige Split, der einen anderen Stapel zeigt als das Quellmaterial und
dessen Ergebnisse in keine Entscheidung eingehen; die
belastbaren Zahlen sind die \code{test}-Zahlen.

\absatz{Grenzen der Ablation.} Drei Vorbehalte bleiben. \emph{Erstens} das
reduzierte Budget: 3000 Bilder und 30 Epochen sind aus der Rechenzeit gewählt,
nicht aus einem Konvergenzkriterium. Es ist für beide Arme gleich und kann
den Vergleich deshalb nicht verzerren, begrenzt aber die absolute Höhe der
Zahlen; ein Effekt, der sich erst nach 60 Epochen zeigte, bliebe unentdeckt.
\emph{Zweitens} je ein Lauf pro Arm: Ohne Wiederholung mit verschiedenen
Trainings-Seeds ist die Streuung des Trainings selbst unbekannt und von der
Wirkung des Datensatzes nicht zu trennen, zumal der Effekt nur auf einem der
beiden Bilder größer ist als die Messunsicherheit. Das ist die gewichtigste
Einschränkung, und sie ist mit Rechenzeit behebbar: Drei Seeds je Arm wären
sechs statt zwei Läufen. \emph{Drittens} ist die Ablation ein \emph{Eingriff
in eine Richtung}: Sie
zeigt die Folgen einer Verschlechterung, nicht den Ertrag eines besseren
Entwurfs derselben Stufen.

\absatz{Nächste Schritte.} Jede der genannten Grenzen hat eine Gegenmaßnahme;
nach erwartetem Nutzen geordnet:

\begin{enumerate}[itemsep=3pt,leftmargin=1.6em]
  \item \textbf{Labelkonvention angleichen.} Eine geänderte Boxschwelle im
  Generator genügt, und die Sensitivitätsrechnung lässt auf \code{test} rund
  \zsensTestGain{} Punkte \mAPseven{} erwarten; das ist der mit geringstem Aufwand erreichbare Gewinn dieser Liste, anzugehen vor jeder weiteren Optimierung.
  \item \textbf{Mehr Bewertungsdaten und mehrere Seeds.} Zehn bis zwanzig
  annotierte Realbilder aus Aufnahmen, aus denen kein Quellmaterial stammt,
  würden die Konfidenzintervalle auf ein Maß bringen, bei
  dem auch kleinere Differenzen entscheidbar wären; drei Trainings-Seeds je Arm
  trennen zusätzlich die Streuung des Trainings von der Wirkung des Datensatzes.
  Zusammen sind das die beiden Punkte, die aus einem Hinweis einen Nachweis
  machen.
  \item \textbf{Das Bündel der Ablation auflösen.} Wie sich der Effekt auf die
  einzelnen Einbettungsmaßnahmen verteilt, ist offen. Je ein Lauf ohne Farb-,
  Rausch- und Unschärfeangleich, ohne weiche Alphakante und ohne JPEG-Durchlauf
  schlüsselt ihn auf, mit dem Sim-to-Real-Abstand als Zielgröße. Schon zwei Arme
  trennen Abkürzung und Domänenlücke (Abschnitt~\ref{sec:ablationC}): Verliert
  der Arm nur mit harter Kante mehr als der nur ohne photometrischen Angleich,
  spricht das für die Abkürzung. Auf dieselbe Weise ließen sich der
  Verdeckungs\emph{grad} und die \emph{Herkunft} der Okkluder prüfen, die
  bisher nur begründet sind.
  \item \textbf{Breiteres Quellmaterial.} Objektausschnitte aus mehreren
  Aufnahmen, darunter Enden mit gefüllter oder verschmutzter Bohrung, die die
  Fehleranalyse als den Fehlertyp ausweist, der nicht auf Objektgröße
  zurückgeht; dazu Clean Plates bei harter Sonne, Gegenlicht und Nässe.
\end{enumerate}

\section{Fazit}

Ein ausschließlich auf synthetischen Bildern trainierter Detektor erreicht auf
zwei echten Aufnahmen \mAPfifty{} von \textbf{\zrefValMap} und
\textbf{\zrefTestMap} bei einem besten $F_1$ von \zrefValFf{} und \zrefTestFf;
auf \code{test}, das weder im Quellmaterial vorkommt noch zur Modellauswahl
dient, zählen von \znPredTest{} Detektionen nur \zfpTest{} als falsch-positiv,
\zfpDoppelTest{} davon als zweite Box auf einer bereits erkannten Stirnfläche. Für
eine Anwendung, für die kein einziges handannotiertes Trainingsbild existiert,
ist das ein tragfähiges Ergebnis.

Methodisch tragen drei Einsichten:

\begin{enumerate}[itemsep=3pt]
  \item \textbf{Bei rein synthetischen Trainingsdaten ist die Pipeline die
  Hypothese.} Jede Entscheidung über eine Transformation ist eine Aussage
  darüber, welche Invarianzen in der Zieldomäne gelten. Die Begründungspflicht
  (Messung, physikalische Randbedingung, Literatur oder ausgewiesene Annahme)
  ist deshalb keine Formalität, sondern das eigentliche Handwerk. Sie hat
  zudem einen Nebeneffekt, den diese Arbeit nutzt: Wer jede Begründung
  ausschreibt, macht sichtbar, welche davon noch nie geprüft wurde.
  \item \textbf{Sorgfalt beim Einfügen ist kein Feinschliff, und keine
  synthetische Kennzahl zeigt an, wenn sie fehlt.} Werden die Maßnahmen gegen
  die Signaturen S1--S5 abgeschaltet, ändert sich \mAPfifty{} auf dem
  Testbild um $\zablCTestGap$\PP, auf dem \emph{eigenen} synthetischen Split
  desselben Arms aber nur um $\zablCSvalGap$\PP; auf dem zweiten Realbild
  bleibt der Unterschied innerhalb der Messunsicherheit. Der Schaden entsteht
  ausschließlich in der Übertragung und ist an den synthetischen Daten
  unsichtbar. Wer eine Copy-Paste-Pipeline an ihren eigenen Bildern misst,
  misst die falsche Größe; nötig ist ein beschrifteter realer Split, und sei
  es ein einziges Bild.
  \item \textbf{Auflösung dominiert, und die Labelkonvention ist Teil des
  Messinstruments.} Zwischen \code{imgsz}~640 und dem Optimum liegen
  \zswValGain\PP{} \mAPfifty, mehr als das Vierfache des Effekts der
  Einbettung. Ein systematischer Größenversatz von rund zwei Pixeln
  zwischen Generatorlabel und Handannotation verschiebt außerdem \mAPseven{}
  auf dem Testbild von \zsensTestMapSOrig{} auf \zsensTestMapSNeu, ohne dass
  sich am Modell etwas ändert. Wer Lokalisierungsmetriken bei 20\px-Objekten berichtet, muss beide
  Größen zuerst kennen.
\end{enumerate}
